\chapter{Cross-run state tracking (menu 49)}
\label{ch:menu49}
\menulabel{49}

\section{Purpose}

Track one electronic state across a sequence of runs of one geometry.
The measure is the density-matrix tracking criterion of Tran, Shea \&
Neuscamman (JCTC 2019) in its post-processing form: every candidate root
of the next run is scored by the Frobenius difference of its one-body
reduced density matrix against the tracked state's density (rotated into
the candidate's orbital basis) plus the energy term $W_0 = (E_t -
E_k)^2$; the smallest $Q = W_0 + D$ continues the lineage.  Use it where
root flipping would otherwise go unnoticed: an SA-size change, a rerun
into a different-shaped solution, or any coarse-to-fine or
method-to-method sequence of one system.

\section{What you need}

One \code{orca\_2json} export per run
(\code{\{"MOCoefficients": true, "1elIntegrals": ["S"], "Densities":
["all"]\}} written next to the gbw; the CASSCF run must have kept its
densities, \code{KeepDens}), and -- optionally -- each export's sibling
\code{.out} for the per-root energies (a missing sibling only costs the
$W_0$ term, which the report then declares).  All runs of the sequence
must share atoms, coordinates and basis (the export's own S matrix,
compared across runs), because the density rotation presupposes one AO
space.

\section{Walkthrough}

The example walks root~0 of the SA(6,6) three-root run into the
four-root run of the same system (the runs differ only in \code{nroots}
-- the tracked state must survive the change of the average):

\lstinputlisting[style=fbkcode, firstline=57, lastline=70,
  title={The menu-49 example (the lineage, the per-candidate table and the
  margin)}]
  {examples/expected/m49_tracking.txt}

The ground state matches at $D = 7.2\times10^{-3}$ while every other
root sits at $0.887$--$0.900$; the same walk from root~1 matches at
$1.26\times10^{-2}$ with the runner-up at $1.84\times10^{-2}$ (margin
$6.7\times10^{-3}$), and from root~2 the margin collapses to
$9\times10^{-16}$: roots 2 and 3 of the four-root run print degenerate,
and the report marks the step ``nearly degenerate in this metric''.

\section{What you get}

The report \code{<first-export>.track.fbk.md}: the run list in tracking
order, the target's label and energy, one table per step (each candidate's
multiplicity, root, energy, $W_0$, $D$ and $Q$, the chosen one marked, and
the margin over the runner-up), and the boundary declarations.  The walk
is deterministic and replays byte-identically.

\section{Conventions, cross-checks and boundaries (measured)}

\begin{readbox}
\begin{itemize}
  \item \textbf{The density channel and its convention.}  The per-root
    densities are the \code{Tdens-CAS.mult.<m>.root.<k>.p} entries of the
    export's Densities sidecar (one per root; the SA(6,6) fixtures ship
    three and four), read with the menu-23 convention $\Gamma = C\,(S\,
    D_{\mathrm{json}}\,S)\,C^{\mathsf T}$; every 1-RDM's trace equals the
    electron count (14 on the N$_2$ fixtures).
  \item \textbf{The cross-run rotation is a measured identity.}  For runs
    of one geometry the S matrices agree to the last digit and
    $M = C_{\mathrm{cand}}\,S\,C_{\mathrm{target}}^{\mathsf T}$ comes out
    orthogonal to $2\times10^{-14}$; the same state then measures
    $D = 7.2\times10^{-3}$ across the average-size change (SA~3 roots
    $\to$ SA~4 roots) and $D = 5.2\times10^{-2}$ across SA $\to$
    single-root (full orbital relaxation), while wrong roots sit at
    $0.887$--$0.900$.
  \item \textbf{Three declared substitutions against the source's
    $Q = W_0 + W_1 + D$.}  The $W_1$ term (the active-to-virtual
    stationarity measure) is not evaluated: no ORCA export carries the
    coupling it needs; for a converged candidate run $W_1$ tends to zero
    (the source's equivalence, $W_1 = 0$ iff the energy is stationary).
    The $1/n_{\mathrm{CAS}}$ scaling of $D$ is not applied (it serves the
    source's portability across active-space sizes; the ranking within
    one sequence is unaffected).  $\omega$ -- the source's energy target
    -- is taken as the tracked state's own energy, updated each step.
  \item \textbf{The margin is the reading.}  The source's own hard case
    (two roots with similarly structured densities) is on record here:
    the printed-degenerate pair shows up as a $10^{-16}$ margin, which
    the report flags rather than resolves; the $10^{-3}$ flagging line is
    presentational, not a criterion.
\end{itemize}
\end{readbox}

\section{Refusals and related sections}

\begin{refusalbox}
\begin{itemize}
  \item Fewer than two runs, a target root index the first run does not
    have (the available labels are listed), runs whose atoms,
    coordinates (tolerance $10^{-6}$~\AA), or AO dimensions differ, an
    export without the S-Matrix, and an export without the per-root
    \code{Tdens-CAS} densities are refused, each with its next step.
  \item Cross-geometry tracking is outside this menu's scope: the
    density rotation presupposes one AO space, and the cross-structure
    orbital mapping (Chapter~\ref{ch:menu17}) remains the route that
    relates two geometries' orbital spaces.  The single-run per-state
    tables are Chapter~\ref{ch:menu31}'s ground; the density channel is
    the menu-23 chain's.
\end{itemize}
\end{refusalbox}
